% paper2_supplement.tex
% Supplemental Material for:
% "Information Loss is Local:
%  The Gravitational Refractive Index as Petz Recovery Fidelity"
% Author: Sheng-Kai Huang
% Compile: pdflatex paper2_supplement.tex

\documentclass[prl,twocolumn,superscriptaddress,longbibliography,floatfix]{revtex4-2}

\usepackage{amsmath}
\usepackage{amssymb}
\usepackage{amsfonts}
\usepackage{hyperref}
\usepackage{booktabs}

% Macros (same as main paper)
\newcommand{\kB}{k_{\mathrm{B}}}
\newcommand{\Tr}{\mathrm{Tr}}
\newcommand{\Petz}{\mathcal{R}}
\newcommand{\chan}{\mathcal{N}}
\newcommand{\tauasym}{\tau}
\newcommand{\ceff}{c_{\mathrm{eff}}}
\newcommand{\rs}{r_{\mathrm{s}}}
\newcommand{\Sgrav}{\Sigma_{\mathrm{grav}}}
\newcommand{\Fbound}{F_{\mathrm{bound}}}

\begin{document}

\title{Supplemental Material for ``Information Loss is Local''}

\author{Sheng-Kai Huang}
\email{akai@fawstudio.com}
\affiliation{Independent Researcher}

\date{\today}

\maketitle

%=======================================================================
\section{Kodama vector in Vaidya spacetime}
\label{sec:kodama_derivation}
%=======================================================================

The ingoing Vaidya metric in advanced
Eddington--Finkelstein coordinates $(v,r,\theta,\phi)$ is
\begin{equation}
  ds^2 = -f\,dv^2 + 2\,dv\,dr + r^2\,d\Omega^2\,,
  \label{eq:vaidya}
\end{equation}
where $f(v,r) \equiv 1 - 2M(v)/r$ and $M(v)$ is the
Bondi mass.
The 2D orbit-space metric on the $(v,r)$ submanifold is
\begin{equation}
  h_{AB}\,dx^A dx^B = -f\,dv^2 + 2\,dv\,dr\,,
  \label{eq:h2d}
\end{equation}
with $\det(h) = -1$.

The Kodama vector~\cite{Kodama1980} is defined as
$K_A = \epsilon_A{}^B\,\partial_B R$, where $R = r$ is
the areal radius and $\epsilon_{AB}$ is the volume form
on the orbit space: $\epsilon_{vr} = \sqrt{-\det h} = 1$.

\emph{Inverse metric.}---From
$h_{AB} = \bigl(\begin{smallmatrix} -f & 1 \\ 1 & 0
\end{smallmatrix}\bigr)$, the inverse is
\begin{equation}
  h^{AB} = \begin{pmatrix} 0 & 1 \\ 1 & f \end{pmatrix}.
  \label{eq:hinv}
\end{equation}

\emph{Mixed Levi-Civita tensor.}---Raising the second index:
\begin{align}
  \epsilon_v{}^v &= \epsilon_{vA}\,h^{Av}
    = \epsilon_{vr}\,h^{rv} = 1\,,
  \label{eq:eps1}\\
  \epsilon_v{}^r &= \epsilon_{vA}\,h^{Ar}
    = \epsilon_{vr}\,h^{rr} = f\,,
  \label{eq:eps2}\\
  \epsilon_r{}^v &= \epsilon_{rA}\,h^{Av}
    = \epsilon_{rv}\,h^{vv} = 0\,,
  \label{eq:eps3}\\
  \epsilon_r{}^r &= \epsilon_{rA}\,h^{Ar}
    = \epsilon_{rv}\,h^{vr} = -1\,.
  \label{eq:eps4}
\end{align}

\emph{Covariant components.}---Since
$\partial_B R = \delta^r_B$:
\begin{align}
  K_v &= \epsilon_v{}^B\,\partial_B R
       = \epsilon_v{}^r = f = 1 - \frac{2M(v)}{r}\,,
  \label{eq:Kv}\\
  K_r &= \epsilon_r{}^B\,\partial_B R
       = \epsilon_r{}^r = -1\,.
  \label{eq:Kr}
\end{align}

\emph{Contravariant components and norm.}---Raising with
$h^{AB}$:
\begin{equation}
  K^v = h^{vA}K_A = K_r = -1\,,\qquad
  K^r = h^{rA}K_A = K_v + f\,K_r = 0\,.
\end{equation}
The norm is
\begin{equation}
  |K|^2 = h_{AB}\,K^A K^B
  = h_{vv}\,(K^v)^2 = -f\,,
  \label{eq:Knorm}
\end{equation}
so
\begin{equation}
  \boxed{\;|K| = \sqrt{1 - \frac{2M(v)}{r}}\;\;}
  \quad\text{for } r > 2M(v)\,.
  \label{eq:Knorm_result}
\end{equation}

The Kodama vector is timelike ($|K|^2 < 0$) for
$r > 2M(v)$, null at $r = 2M(v)$, and spacelike for
$r < 2M(v)$.
Crucially, $|K|$ depends only on the \emph{current} mass
$M(v)$ and radius $r$---it is a local quantity.

The temporal asymmetry parameter is therefore
\begin{equation}
  \tauasym_K(r,v) = 1 - |K| = 1 - \sqrt{1 - \frac{2M(v)}{r}}\,,
  \label{eq:tauK}
\end{equation}
reaching $\tauasym_K = 1$ at $r = 2M(v)$, the
\emph{apparent horizon}.

%=======================================================================
\section{Event horizon in Vaidya spacetime}
\label{sec:EH}
%=======================================================================

The event horizon (EH) is the outermost outgoing null
geodesic that just fails to escape to $r \to \infty$.
Outgoing radial null geodesics in the Vaidya
metric~\eqref{eq:vaidya} satisfy
\begin{equation}
  \frac{dr}{dv} = \frac{1}{2}\!\left(1
  - \frac{2M(v)}{r}\right).
  \label{eq:null_ode}
\end{equation}
The EH is the solution $r_{\text{EH}}(v)$
of~\eqref{eq:null_ode} with boundary condition
$r_{\text{EH}}(v) \to 2M_{\text{final}}$ as
$v \to \infty$.
This must be integrated \emph{backward} from
$v = \infty$---a manifestly teleological procedure.

\subsection{Step-function collapse}

For $M(v) = M_0\,\theta(v)$:
\begin{itemize}
\item $v > 0$ (Schwarzschild):
  $dr/dv = \frac{1}{2}(1 - 2M_0/r)$, static EH at
  $r = 2M_0$.
\item $v < 0$ (Minkowski): $dr/dv = 1/2$, so
  $r = v/2 + C$.  Matching at $v = 0$ gives
  $r_{\text{EH}}(v) = v/2 + 2M_0$.
\end{itemize}
The EH originates at $r = 0$ when $v = -4M_0$---in flat
space, \emph{before} the shell arrives.  The apparent
horizon appears only at $v = 0$ at $r = 2M_0$.

\subsection{Self-similar accretion: $M(v) = \lambda v$}
\label{sec:selfsimilar}

The null geodesic equation~\eqref{eq:null_ode} admits a
self-similar ansatz $r(v) = C\,v$:
\begin{equation}
  C = \frac{1}{2}\!\left(1 - \frac{2\lambda}{C}\right)
  \quad\Longrightarrow\quad
  2C^2 - C + 2\lambda = 0\,.
  \label{eq:selfsimilar}
\end{equation}
For $\lambda < 1/16$, two real roots:
\begin{equation}
  C_{\pm} = \frac{1 \pm \sqrt{1 - 16\lambda}}{4}\,.
  \label{eq:Cpm}
\end{equation}
The event horizon is $r_{\text{EH}}(v) = C_{-}\,v$.
The apparent horizon is $r_{\text{AH}}(v) = 2\lambda\,v$.

\emph{Proof that $C_{-} > 2\lambda$} (i.e.,
$r_{\text{EH}} > r_{\text{AH}}$).---Define
$g(\lambda) \equiv C_{-} - 2\lambda
= \frac{1}{4}(1 - \sqrt{1-16\lambda}) - 2\lambda$.
At $\lambda = 0$: $g(0) = 0$.
The derivative
\begin{equation}
  g'(\lambda) = \frac{2}{\sqrt{1-16\lambda}} - 2
  = 2\!\left(\frac{1}{\sqrt{1-16\lambda}} - 1\right)
  > 0
\end{equation}
for all $\lambda > 0$.  Since $g(0) = 0$ and
$g' > 0$, we have $g(\lambda) > 0$ for
$\lambda \in (0, 1/16)$.  \hfill$\square$

For small $\lambda$:
\begin{equation}
  \frac{r_{\text{EH}}}{r_{\text{AH}}}
  = \frac{C_{-}}{2\lambda}
  = 1 + 4\lambda + O(\lambda^2)\,,
\end{equation}
so the fractional excess grows linearly with accretion
rate.

\subsection{$\tauasym_K$ at the event horizon}

At $r = r_{\text{EH}} = C_{-}\,v$ with
$M(v) = \lambda v$:
\begin{equation}
  \frac{2M(v)}{r_{\text{EH}}}
  = \frac{2\lambda}{C_{-}}\,.
  \label{eq:ratio}
\end{equation}
This is a constant (independent of $v$!).
The $\tauasym$ deficit at the EH is
\begin{equation}
  \tauasym_K(r_{\text{EH}}) = 1
  - \sqrt{1 - \frac{2\lambda}{C_{-}}} < 1\,.
  \label{eq:tau_EH}
\end{equation}

Table~I of the main text is computed from
Eqs.~\eqref{eq:Cpm}--\eqref{eq:tau_EH}.
For example, at $\lambda = 0.03$:
$C_{-} = (1-\sqrt{0.52})/4 = 0.0697$,
$2\lambda/C_{-} = 0.861$,
$\tauasym_K = 1 - \sqrt{0.139} = 0.627$.

\subsection{Smooth accretion: $M(v) = M_0(1-e^{-v/v_0})$}

At late times ($v \gg v_0$), linearization of the null
ODE around $r = 2M_0$ gives the particular solution
\begin{equation}
  r_{\text{EH}}(v) \approx 2M_0
  - \frac{2M_0 v_0}{4M_0 + v_0}\,e^{-v/v_0}\,.
  \label{eq:rEH_smooth}
\end{equation}
The difference
\begin{equation}
  \Delta r \equiv r_{\text{EH}} - r_{\text{AH}}
  = \frac{8M_0^2}{4M_0 + v_0}\,e^{-v/v_0} > 0
  \label{eq:Delta_r}
\end{equation}
confirms $r_{\text{EH}} > r_{\text{AH}}$ during
accretion, with the gap proportional to the accretion
rate $dM/dv = (M_0/v_0)\,e^{-v/v_0}$.

%=======================================================================
\section{Evaporating black hole}
\label{sec:evaporation}
%=======================================================================

For an evaporating black hole with linearly decreasing
mass $M(v) = M_0(1 - v/v_{\text{evap}})$, the
situation reverses~\cite{Hiscock1981}:
\begin{equation}
  r_{\text{AH}}(v) = 2M(v) > r_{\text{EH}}(v)\,.
\end{equation}
The EH is \emph{inside} the AH because it ``knows''
the mass will decrease.  The Kodama vector is spacelike
at $r_{\text{EH}}$ (since $r_{\text{EH}} < 2M(v)$,
the argument of the square root in
Eq.~\eqref{eq:Knorm_result} is negative), confirming
that the EH is inside the trapped region.

The sign of $\Delta r = r_{\text{EH}} - r_{\text{AH}}$
encodes the direction of mass change:
\begin{center}
\begin{tabular}{lcc}
  \textbf{Scenario} & $dM/dv$ & $\Delta r$ \\
  \colrule
  Accretion & $> 0$ & $> 0$ (EH outside AH) \\
  Static    & $= 0$ & $= 0$ (EH $=$ AH) \\
  Evaporation & $< 0$ & $< 0$ (EH inside AH) \\
\end{tabular}
\end{center}

In all cases, $\tauasym_K = 1$ at the AH, never at the EH.
The event horizon requires future knowledge;
the apparent horizon (and $\tauasym$) does not.

%=======================================================================
\section{Connection to quantum extremal surfaces}
\label{sec:QES}
%=======================================================================

The quantum extremal surface
(QES)~\cite{Engelhardt2015,Penington2020}
extremizes the generalized entropy
\begin{equation}
  S_{\text{gen}} = \frac{A}{4G\hbar} + S_{\text{bulk}}
\end{equation}
on a time-slice.  In the classical limit
($\hbar \to 0$), $S_{\text{bulk}} \to 0$ and the QES
reduces to the classical extremal surface, which for
spherically symmetric spacetimes is the apparent horizon
(the outermost MOTS where $\partial_r A = 0$
$\Leftrightarrow$ $\theta_+ = 0$).

Including quantum corrections:
\begin{equation}
  r_{\text{QES}} = r_{\text{AH}}
  + O\!\left(\frac{\hbar}{M^2}\right),
\end{equation}
where the correction comes from $\partial_r S_{\text{bulk}}$.
During evaporation, $\partial_r S_{\text{bulk}} > 0$
(bulk entropy increases outward due to Hawking pairs),
pushing the QES slightly inward:
$r_{\text{EH}} < r_{\text{QES}} \lesssim r_{\text{AH}}$.

The hierarchy is:
\begin{equation}
  r_{\text{EH}} \;<\; r_{\text{QES}}
  \;\lesssim\; r_{\text{AH}}
  \qquad\text{(during evaporation)}.
\end{equation}

In the $\tauasym$ framework:
\begin{itemize}
\item $\tauasym_K = 1$ at $r_{\text{AH}}$ (classical);
\item $\tauasym_K^{\text{quantum}} = 1$ at
  $r_{\text{QES}}$ (quantum-corrected);
\item $\tauasym_K \neq 1$ at $r_{\text{EH}}$ in either case.
\end{itemize}

The $\tauasym$ framework and the QES prescription agree:
the operationally relevant boundary for information
dynamics is the (quantum-corrected) apparent horizon,
not the event horizon.

%=======================================================================
\section{Echo time delay}
\label{sec:echo}
%=======================================================================

The tortoise coordinate for the exponential metric
$g_{00} = -e^{-\rs/r}$,
$g_{rr} = e^{+\rs/r}$ is
\begin{equation}
  r^{*} = \int \frac{\sqrt{g_{rr}}}{\sqrt{-g_{00}}}\,dr
  = \int e^{\rs/r}\,dr\,.
  \label{eq:tortoise}
\end{equation}
Since $e^{\rs/r}$ is bounded for $r > 0$ (it equals $e$
at $r = \rs$ and approaches 1 as $r \to \infty$), the
tortoise coordinate remains finite.
This contrasts with Schwarzschild, where
$r^{*} \to -\infty$ as $r \to \rs$.

The round-trip light travel time from the photon sphere
($r = \rs$ in isotropic coordinates) through the wormhole
throat and back is
\begin{equation}
  \Delta t_{\text{echo}} = 2 \int_{0}^{\rs}
  e^{\rs/r}\,dr \approx 4.2\,\frac{\rs}{c}\,,
  \label{eq:echo_delay}
\end{equation}
where the numerical coefficient is evaluated by
substituting $u = \rs/r$:
$\int_0^{\rs} e^{\rs/r}\,dr
= \rs \int_1^{\infty} u^{-2}\,e^u\,du$.
The integral converges and equals approximately
$2.1\,\rs$, giving a round trip of $\approx 4.2\,\rs/c$.

For a $60\,M_\odot$ remnant:
$\rs = 2GM/c^2 = 177$~km,
$\Delta t \approx 4.2 \times 177/c \approx 2.5$~ms,
at a frequency $\sim 400$~Hz---squarely in the LIGO band.

%=======================================================================
% BIBLIOGRAPHY
%=======================================================================

\begin{thebibliography}{10}

\bibitem{Kodama1980}
H.~Kodama,
Prog.\ Theor.\ Phys.\ \textbf{63}, 1217 (1980).

\bibitem{Hayward1994}
S.~A.~Hayward,
Phys.\ Rev.\ D \textbf{49}, 6467 (1994).

\bibitem{AbreuVisser2010}
G.~Abreu and M.~Visser,
Phys.\ Rev.\ D \textbf{82}, 044027 (2010).

\bibitem{Hiscock1981}
W.~A.~Hiscock,
Phys.\ Rev.\ D \textbf{23}, 2823 (1981).

\bibitem{Engelhardt2015}
N.~Engelhardt and A.~C.~Wall,
J.\ High Energy Phys.\ \textbf{2015}(01), 073 (2015).

\bibitem{Penington2020}
G.~Penington,
J.\ High Energy Phys.\ \textbf{2020}(09), 002 (2020).

\end{thebibliography}

\end{document}
